\note{sa page 6 ; la rédaction du n° 6 reprend ici}

\[ (6.1)_{i} \qquad L^{n-i} \,|\, M^{i} \colon M^{i} \to M^{2n-i} \ \text{ est un isom.} \]

Ceci implique déjà $(4.1)_{i}$ et $(4.2)_{i}$ donc $(3.1)_{i}$ et $(3.2)_{i}$,
\add{donc l'unicité des $\pi_{i}$ en termes de la structure graduée de
$\mathcal{E}^{*}$.} Supposons de plus

$(6.2)_{i}$ \quad L'inverse de l'isom.\ (6.1) \uncertain{est} induit par un élément de
$\mathcal{E}^{-(2n-2i)}$.

\note{Après « (6.1) », une surcharge biffée, où se lit encore un
$\mathcal{E}$.}

\struck{Ceci} L'ensemble des deux conditions (6.1) et (6.2) équivaut [en termes
intrinsèques : $\mathcal{E}^{*}$ muni de $L \in \mathcal{E}^{2}$ \ill{}
\uncertain{supposant} déjà les $\pi_{i} \in \mathcal{E}^{*}$] à

\[ (6.3)_{i} \qquad \begin{cases} \exists\, w_{i} \in \mathcal{E}^{-(2n-2i)} \ \text{tel qu'on ait} \\[2pt] (w_{i} L^{n-i} - 1)\, \pi_{i} = 0 , \qquad (L^{n-i} w_{i} - 1)\, \pi_{2n-i} = 0 \end{cases} \]

Cette condition \ill{} : elle est indépendante de la donnée particulière d'un
$M^{*}$. \struck{\ill{} donc (3.1) et (3.2)} Elle \ill{} $(4.4)_{i}$, et
(comme on a vu aux n° 4, 5) implique $\pi_{i} \in \mathcal{E}$ et l'unicité des
$\pi_{i}$, en termes des conditions (4.1) et (4.2).

Revenons à $M^{*}$, et $L$ satisfaisant (6.1). Nous avons vu comment
\struck{\ill{}} cette condition implique la décomposition de $M^{*}$ en termes de
morceaux primitifs :

\[ (6.4) \qquad \begin{cases} M^{i} \simeq P^{i} + L P^{i-2} + L^{2} P^{i-4} + \ldots & (i \leqslant n) \\[2pt] M^{2n-i} \simeq L^{n-i} P^{i} + L^{n-i+1} P^{i-2} + \ldots \end{cases} \]

On désigne par $\pi_{i,\alpha}$ la projection de $M^{i}$ sur
$L^{\alpha} P^{i-2\alpha}$, par $\pi_{2n-i,\alpha}$ celle de $M^{2n-i}$ sur
$L^{n-i+\alpha} P^{i-2\alpha}$.

